\chapter{结论}
%%====================================================================

本文以比特币（Bitcoin）开源社区为实证场景，围绕”信息如何驱动开源参与”这个核心问题，从两个互补维度展开考察。研究一聚焦外部组织的捐赠信息披露对个体开发者活动的影响，研究二探究社区讨论情感对集体层面代码活动与讨论活动的动态影响。两项研究在微观个体与宏观社区、横截面事件与纵向动态过程、捐赠信息与情感信息等多个维度上形成互补，共同构建了围绕开源社区信息环境的双层分析框架，为理解信息因素如何塑造开源生态可持续性提供了新的视角。

\section{主要研究结论}

\subsection{捐赠信息披露的激励效应}

研究一以Bitcoin Brink基金会2022年至2023年间的捐赠信息披露事件为准自然实验，采用面板个体固定效应模型与两阶段最小二乘法，基于169,785个开发者—周面板观测（1,610名开发者），考察了不同类型捐赠信息披露对开发者代码活动（提交、拉取请求）和审查活动（代码审查、评论）的影响。

核心发现表明，非定向捐赠披露（仅公告有捐赠活动而不点名受益开发者）对开发者代码活动与审查活动均产生显著正向效应（H1、H2成立）。该结果可由期望理论的核心逻辑加以解释：社区成员获知有捐赠资金进入生态系统时，其对”努力→绩效→奖励”链条的主观期望随之提升——即便尚无法确知自己是否会成为受益者，信息的存在本身已足以激活参与动机。这一”预期性激励（Anticipatory Incentive）”机制揭示了信息的独立价值：它先于实际财务转移而发生作用，先于奖励兑现而改变行为。

定向捐赠披露（公告具体获得资助的开发者名单）同样正向影响了审查活动（H2成立）。与非定向披露的弥散性激励不同，定向披露通过强化”绩效—奖励”因果链的工具性感知，使奖励可得性由模糊变为清晰，为开发者提供了可见的努力方向与明确的参照标准。

然而，研究一同时揭示了该激励机制内部的矛盾性。当定向披露频率增大时，非定向披露的激励效应显著衰减（H3在2SLS模型中成立）。该”期望失验（Expectation Disconfirmation）”效应表明，定向披露实际上引入了一种隐性的比较压力——在增强部分开发者激励的同时，削弱了另一部分开发者对奖励可得性的主观期望。该调节效应在FE模型中方向为负但统计不显著，经2SLS修正内生性偏误后方达显著水平，说明内生性问题会掩盖定向披露的负向调节作用。基金会的披露策略并非孤立地作用于当期行为，而是通过改变开发者心理预期在社区内积累长期效应；定向与非定向披露的频率组合因此需要审慎权衡。

工具变量法（以Google搜索趋势指数的四周移动平均值及其与定向披露的交互项为工具变量）进一步确认，上述激励效应独立于市场价格波动带来的财富效应，且在剔除实际获得捐款的开发者样本后依然稳健。三组替代性度量的稳健性检验（虚拟变量加时间窗口、连续值当期度量、虚拟变量当期度量）及两组机制分析（比特币项目投入程度与开源社区活跃程度的异质性）均支持主要结论。这意味着研究一所识别的是纯粹的信息披露效应而非货币激励效应，在方法论上为”信息即激励”命题提供了较为严格的因果证据。

\subsection{社区情感的动态影响机制}

研究二将分析层次从个体转至社区，采用向量自回归（VAR(2)）模型对比特币GitHub仓库2020年至2023年的周度时间序列数据进行动态分析，探究社区讨论情感与开发者行为活动之间的互动因果关系。控制市场回报、捐赠金额与资助授予等外生变量后，主要发现如下。

社区讨论中情感极性的下降（即负面情感增加）在滞后一期对代码活动产生了显著正向影响（H4成立）。该看似反直觉的效应可由”情感即信息（Affect-as-Information）”理论加以解释\cite{schwarz1983moods}：社区中弥漫的负面情感作为集体信号，传递了当前状态与理想状态之间存在差距的信息，开发者对该”问题信号”的回应是加速代码产出以修复已知缺陷或应对技术压力。在技术型开源社区中，情感并非仅仅是情绪的排泄，它承担着重要的信息传递功能。

负面情感在滞后一期亦对讨论活动产生正向影响（H5成立），与H4方向一致。情感信号的激活效应并不局限于代码生产领域，同样唤起了开发者对话题讨论的参与意愿。社区处于情绪低谷时，开发者倾向于以更多讨论来回应——既包含对技术问题的深入分析，也包含对社区走向的集体关切，讨论活动兼具问题诊断与情绪疏导的双重功能。

代码活动的增加在滞后两期对讨论活动产生显著负向影响（H6成立）。该”替代效应”与目标—手段替代（Goal-Means Substitution）理论高度吻合\cite{kruglanski2002theory}：开发者已通过代码提交实现对社区问题的应对目标后，借助讨论参与达成同一目标的需求随之降低，两种行为在时间和精力上形成竞争关系。这揭示了开源社区行为活动中的内部制衡机制——代码生产并非简单叠加于讨论之上，二者之间存在动态的资源竞争与功能替代。

格兰杰因果检验、脉冲响应函数（IRF）分析与预测误差方差分解（FEVD）进一步确认了上述三条影响路径的统计显著性与经济实质，并将其整合为”情感→代码/讨论→讨论”的完整动态机制。替换情感变量操作化方式（纯负向情感指标）、替换讨论变量、调整乔列斯基排序等多种稳健性检验方案均支持核心发现的稳定性；计划中的传统词典方法（VADER）对照检验将进一步评估结论对情感分析工具选择的不敏感性。

\subsection{两项研究的综合结论}

综合两项研究的发现，可以提炼出一个更具一般性的命题：在开源社区中，信息而非金钱，才是影响开发者参与行为的直接驱动因素。研究一表明，捐赠行为的信息披露——甚至是捐赠可能性的模糊预期——足以产生显著激励效应，无需实际财务转移的发生。研究二表明，情感作为一种非正式集体信息，能够以可预测的方向和时序影响后续行为活动，代码生产与讨论参与均对情感信号作出有规律的反应。二者共同构成”信息驱动参与（Information-Driven Participation）”这一理论命题在不同层次、不同类型信息上的双侧证据。

该命题对理解开源软件可持续性具有重要意义。资金支持固然是维持开源社区长期活力的必要条件，但信息披露机制的设计（以最大化激励效应）、社区情感信息环境的感知与管理，同样是不可忽视的治理维度。开源生态的可持续性在相当程度上取决于信息流通的质量与设计——这是两项研究共同指向的核心结论。

\section{研究展望}

本文在比特币开源社区中取得了若干具有理论与实践价值的发现，但在外部效度、心理机制的直接验证、情感分析的精细化及两项研究的整合等方面仍有深化空间。以下从四个方向提出拓展路径。

\subsection{研究场景的拓展与跨社区比较}

本文以比特币社区为单一研究场景。该社区在规模、历史和开源贡献制度方面具有典型性，但结论的外部效度仍待系统检验。未来研究应将范围扩展至以太坊（Ethereum）、Litecoin等其他加密货币开源社区，以及更广泛的非加密货币开源项目（如Linux内核、NumPy、scikit-learn等）。跨社区比较有助于识别哪些发现具有普遍性（如情感的信息传递功能），哪些依赖于特定的社区制度背景（如Brink基金会的定向资助模式）。

尤需指出的是，研究一中”期望失验”效应在很大程度上依赖于Bitcoin Brink特有的”定向/非定向二元披露结构”。在采用不同资助披露模式的社区中（如GitHub Sponsors的个人订阅模式、Open Collective的众筹模式），类似效应是否存在，需通过专门设计的准自然实验加以检验。跨社区比较的积累将使双层信息环境框架在更广泛的开源生态中得到检验与精炼。

\subsection{心理机制的直接验证与混合方法}

两项研究均以间接推断的方式检验心理机制：研究一通过行为数据的差异推断期望理论的激活，研究二通过情感与行为的时序关系推断情感即信息效应。这种方法论取向在外部效度上具有优势，但无法直接触及开发者的主观认知状态。

未来研究可结合调查问卷或深度访谈，直接测量开发者对捐赠信息披露事件的期望感知（期望值、工具性感知、效价）及其对社区情感氛围的主观体验。”行为数据+主观报告”的混合方法将有助于厘清理论机制的准确性，并区分竞争性解释。例如，研究二中负面情感对代码活动的正向效应，究竟经由认知层面的”问题识别与应对”路径（情感即信息）还是情绪传染的”群体压力”路径发挥作用？主客观数据的结合可使该理论分野得到更直接的检验。

\subsection{情感测量的深化与多维分析}

研究二以大型语言模型（claude-sonnet-4.5）对 Issue 评论进行情感打分，已收集愤怒、厌恶、恐惧、快乐、悲伤、惊讶及期望失验七个细粒度情感维度的评分数据（详见附录D）。受限于研究框架，本文 VAR 模型仅使用了综合情感极性（整体正负向），七类细粒度情感数据尚待充分挖掘。未来研究可沿以下方向推进。一是利用已有细粒度数据，分析愤怒、恐惧、悲伤、期望失验等具体情感类型对代码活动与讨论活动的差异化影响——“愤怒”是否比”沮丧”更能激发即时代码行动？”期望失验”是否与 H3 中个体期望失验机制在社区层面产生共振？二是对 LLM 评分结果开展人工标注基准测试（annotation benchmark），在比特币 Issue 评论样本上评估 claude-sonnet-4.5 打分与人工标注的一致性（Cohen's $\kappa$），进一步确认方法效度。三是引入情感来源的区分维度，将核心开发者（committers）与外围参与者（用户、贡献者）的情感表达分别建模，识别不同来源情感信号的影响异质性及其社区传播路径。

\subsection{两项研究的整合分析框架}

本文将两项研究分别置于个体层面和社区层面加以考察。第六章虽从理论层面提出了”双层信息环境”框架并讨论了跨层次反馈假说，但尚未在实证层面直接检验两类信息效应的交互作用。未来研究可尝试构建整合性分析框架，同时纳入捐赠信息披露事件与社区情感变量，考察二者的独立效应、交互效应及跨层次反馈机制。

具体而言：捐赠信息披露对个体行为的激励效应，是否因社区情感氛围的差异而产生调节？在负面情感弥漫的社区背景下，非定向捐赠披露的”预期性激励”会被削弱还是放大？反过来，个体层面激励机制的激活，是否会通过行为外溢影响社区整体情感状态，从而形成从信息披露到情感再到行为的完整反馈回路？这些跨层次动态过程是理解开源社区信息生态的重要议题，也是双层信息环境框架最具潜力的发展方向。方法层面，未来可考虑引入多层次模型（HLM）、门限VAR（Threshold VAR）或时变参数VAR（TVP-VAR）模型，以更完整地刻画信息—情感—行为三者之间的非线性与时变特征。


%---------------------------------------------------------------------
%   参考文献
%---------------------------------------------------------------------

\printbibliography

%---------------------------------------------------------------------
%   致谢
%---------------------------------------------------------------------

\begin{acknowledgement}
感谢导师的悉心指导与耐心支持。
感谢图书情报学院各位老师在研究方法和学术写作方面的宝贵建议。
感谢 Bitcoin 开源社区的全体开发者，正是他们的开放协作精神使本研究成为可能。
感谢 Bitcoin Brink 基金会公开其运营数据，为本研究提供了独特的研究场景。
感谢同门师兄弟姐妹在讨论中给予的灵感和鼓励。
\end{acknowledgement}

%---------------------------------------------------------------------
%   附录
%---------------------------------------------------------------------

